\documentclass[11pt,a4paper]{article}

\usepackage[utf8]{inputenc}
\usepackage[T1]{fontenc}
\usepackage{mathptmx} % Font Times New Roman
\usepackage{amsmath,amssymb,amsfonts}
\usepackage{graphicx}
\usepackage{booktabs}
\usepackage{geometry}
\usepackage{xcolor}
\usepackage{listings}
\usepackage{array}
\usepackage{hyperref}

\geometry{
    top=2.5cm,
    bottom=2.5cm,
    left=2.5cm,
    right=2.5cm
}

\graphicspath{{images/}{./}}

% Helper command agar kompilasi tidak error jika gambar belum diupload di Overleaf
\newcommand{\insertsafeimg}[2]{%
    \IfFileExists{images/#1}{\includegraphics[#2]{images/#1}}{%
        \IfFileExists{#1}{\includegraphics[#2]{#1}}{%
            \fbox{\parbox{0.85\textwidth}{\centering\vspace{0.8cm}\textbf{[Gambar: #1]}\\\small\textit{(Silakan upload #1 ke direktori images/ di Overleaf jika belum ada)}\vspace{0.8cm}}}%
        }%
    }%
}

% -----------------------------------------------------------------------------
% Konfigurasi Code Listing (C & Terminal Output)
% -----------------------------------------------------------------------------
\definecolor{codebg}{rgb}{0.97,0.97,0.98}
\definecolor{commentgreen}{rgb}{0.13,0.54,0.13}
\definecolor{keywordblue}{rgb}{0.0,0.2,0.8}
\definecolor{stringpurple}{rgb}{0.58,0.0,0.82}
\definecolor{linenumgray}{rgb}{0.5,0.5,0.5}
\definecolor{bordergray}{rgb}{0.8,0.8,0.8}

\lstdefinestyle{ccode}{
    backgroundcolor=\color{codebg},
    commentstyle=\color{commentgreen}\textit,
    keywordstyle=\color{keywordblue}\bfseries,
    numberstyle=\tiny\color{linenumgray},
    stringstyle=\color{stringpurple},
    basicstyle=\ttfamily\footnotesize,
    breakatwhitespace=false,
    breaklines=true,
    captionpos=b,
    keepspaces=true,
    numbers=left,
    numbersep=7pt,
    showspaces=false,
    showstringspaces=false,
    showtabs=false,
    tabsize=4,
    frame=single,
    rulecolor=\color{bordergray}
}

\lstdefinestyle{cmdoutput}{
    backgroundcolor=\color{gray!6},
    basicstyle=\ttfamily\scriptsize,
    breaklines=true,
    frame=single,
    rulecolor=\color{bordergray},
    numbers=none
}

\hypersetup{
    colorlinks=true,
    linkcolor=blue!70!black,
    urlcolor=blue!70!black
}

\begin{document}

% =============================================================================
% HALAMAN COVER (Sesuai Format Akademik FTMM Universitas Airlangga)
% =============================================================================
\begin{titlepage}
    \centering
    {\fontsize{16pt}{20pt}\selectfont\textbf{LAPORAN TUGAS PRAKTIKUM}}\\[0.3cm]
    {\fontsize{14pt}{18pt}\selectfont\textbf{(SIA103) STRUKTUR DATA}}\\[0.25cm]
    {\fontsize{12pt}{15pt}\selectfont\textbf{Topik : Queue (Antrian)}}\\[0.15cm]
    {\fontsize{12pt}{15pt}\selectfont\textbf{Pertemuan ke-5 (TM 5)}}\\[1.2cm]

    \insertsafeimg{logo_unair.png}{width=0.35\textwidth}\\[1.2cm]

    {\fontsize{11pt}{14pt}\selectfont\textbf{Dosen Pengampu:}}\\[0.15cm]
    Purbandini, S.Si., M.Kom.\\
    Rodik Wahyu Indrawan, S.ST., M.Tr.T.\\[1.0cm]

    {\fontsize{11pt}{14pt}\selectfont\textbf{Oleh:}}\\[0.25cm]
    {\fontsize{12pt}{15pt}\selectfont\textbf{Hellyos Ageng Haqiqie}}\\[0.1cm]
    {\fontsize{11pt}{14pt}\selectfont\textbf{(163251001)}}\\[1.2cm]

    {\fontsize{12pt}{16pt}\selectfont\textbf{Program Studi Teknik Robotika dan Kecerdasan Buatan}}\\[0.1cm]
    {\fontsize{12pt}{16pt}\selectfont\textbf{Fakultas Teknologi Maju dan Multidisiplin}}\\[0.1cm]
    {\fontsize{12pt}{16pt}\selectfont\textbf{Universitas Airlangga}}\\[0.1cm]
    {\fontsize{12pt}{16pt}\selectfont\textbf{2026}}
\end{titlepage}

\newpage

% =============================================================================
% SOAL 1: PROGRAM MENU QUEUE LINEAR ARRAY
% =============================================================================
\begin{center}
    {\Large\textbf{LATIHAN SOAL 1}}
\end{center}

\vspace{0.2cm}
\noindent\textbf{Soal:}
\begin{quote}
    \textit{"Buatlah program dalam bentuk menu untuk mengimplementasikan antrian. Menu tsb berisi memasukkan data, menghapus data, menampilkan data dan keluar."}
\end{quote}

\vspace{0.3cm}
\noindent\textbf{1. Kode Program (\texttt{soal1\_queue\_linear.c}):}

\begin{lstlisting}[style=ccode, language=C]
#include <stdio.h>
#include <stdlib.h>

#define MAX 8

typedef struct {
    int data[MAX];
    int head;
    int tail;
} Queue;

Queue antrian;

/* Inisialisasi antrian kosong */
void Create(void) {
    antrian.head = -1;
    antrian.tail = -1;
}

/* Memeriksa apakah antrian kosong */
int IsEmpty(void) {
    if (antrian.tail == -1) {
        return 1;
    }
    return 0;
}

/* Memeriksa apakah antrian penuh */
int IsFull(void) {
    if (antrian.tail >= MAX - 1) {
        return 1;
    }
    return 0;
}

/* Menambahkan elemen ke dalam antrian */
void Enqueue(int data) {
    if (IsFull() == 1) {
        printf("[PERINGATAN] Antrian Penuh (Overflow)! Tidak dapat menambahkan data %d.\n", data);
        return;
    }

    if (IsEmpty() == 1) {
        antrian.head = 0;
        antrian.tail = 0;
        antrian.data[antrian.tail] = data;
    } else {
        antrian.tail++;
        antrian.data[antrian.tail] = data;
    }
    printf("[SUKSES] Elemen %d berhasil masuk ke antrian (posisi indeks: %d).\n", data, antrian.tail);
}

/* Mengambil dan menghapus elemen terdepan dari antrian */
int Dequeue(void) {
    if (IsEmpty() == 1) {
        printf("[PERINGATAN] Antrian Kosong (Underflow)! Tidak ada elemen yang dapat dikeluarkan.\n");
        return -1;
    }

    int elemenKeluar = antrian.data[antrian.head];

    /* Geser seluruh elemen ke arah depan menggunakan looping */
    for (int i = antrian.head; i < antrian.tail; i++) {
        antrian.data[i] = antrian.data[i + 1];
    }
    antrian.tail--;

    /* Jika setelah digeser elemen menjadi habis, kembalikan ke kondisi awal */
    if (antrian.tail < 0) {
        antrian.head = -1;
        antrian.tail = -1;
    }

    printf("[SUKSES] Elemen %d berhasil keluar dari antrian.\n", elemenKeluar);
    return elemenKeluar;
}

/* Mengosongkan antrian */
void Clear(void) {
    antrian.head = -1;
    antrian.tail = -1;
    printf("[SUKSES] Seluruh isi antrian berhasil dibersihkan (Clear).\n");
}

/* Menampilkan seluruh elemen dalam antrian */
void Tampil(void) {
    if (IsEmpty() == 1) {
        printf("[INFO] Antrian kosong! Tidak ada data untuk ditampilkan.\n");
        return;
    }

    printf("\n=== STATUS ELEMEN ANTRIAN (LINEAR ARRAY) ===\n");
    printf("Kapasitas Maksimum : %d elemen\n", MAX);
    printf("Jumlah Elemen Saat Ini: %d elemen\n", antrian.tail - antrian.head + 1);
    printf("Indeks Head (Depan)   : %d (Nilai: %d)\n", antrian.head, antrian.data[antrian.head]);
    printf("Indeks Tail (Belakang): %d (Nilai: %d)\n", antrian.tail, antrian.data[antrian.tail]);
    printf("Daftar Antrian: [FRONT] ");
    for (int i = antrian.head; i <= antrian.tail; i++) {
        printf("%d", antrian.data[i]);
        if (i < antrian.tail) {
            printf(" -> ");
        }
    }
    printf(" [REAR]\n");

    /* Tampilan visualisasi slot memori array */
    printf("Slot Memori Array:\n");
    printf("+-------");
    for (int i = 0; i < MAX; i++) printf("+-------");
    printf("+\n| Indeks");
    for (int i = 0; i < MAX; i++) printf("|  [%d]  ", i);
    printf("|\n+-------");
    for (int i = 0; i < MAX; i++) printf("+-------");
    printf("+\n| Nilai ");
    for (int i = 0; i < MAX; i++) {
        if (i >= antrian.head && i <= antrian.tail) {
            printf("| %5d ", antrian.data[i]);
        } else {
            printf("|   -   ");
        }
    }
    printf("|\n+-------");
    for (int i = 0; i < MAX; i++) printf("+-------");
    printf("+\n\n");
}

int main(void) {
    int pilihan, nilai;
    Create();

    printf("========================================================\n");
    printf("   PROGRAM IMPLEMENTASI QUEUE LINEAR ARRAY (SOAL 1)    \n");
    printf("   Nama : Hellyos Ageng Haqiqie (NIM: 163251001)       \n");
    printf("   Prodi: S1 Teknik Robotika dan Kecerdasan Buatan      \n");
    printf("========================================================\n");

    do {
        printf("\n========================================\n");
        printf("           MENU PILIHAN UTAMA           \n");
        printf("========================================\n");
        printf(" 1. Masukkan Data ke Antrian (Enqueue) \n");
        printf(" 2. Hapus Data dari Antrian (Dequeue)  \n");
        printf(" 3. Tampilkan Data Antrian (Tampil)    \n");
        printf(" 4. Bersihkan Seluruh Antrian (Clear)  \n");
        printf(" 5. Keluar dari Program                \n");
        printf("========================================\n");
        printf("Pilihan Anda [1-5]: ");
        if (scanf("%d", &pilihan) != 1) {
            printf("[ERROR] Input tidak valid! Harap masukkan angka.\n");
            while (getchar() != '\n');
            continue;
        }

        switch (pilihan) {
            case 1:
                printf("Masukkan nilai bilangan bulat yang akan ditambahkan: ");
                if (scanf("%d", &nilai) == 1) {
                    Enqueue(nilai);
                } else {
                    printf("[ERROR] Input harus berupa bilangan bulat!\n");
                    while (getchar() != '\n');
                }
                break;
            case 2:
                Dequeue();
                break;
            case 3:
                Tampil();
                break;
            case 4:
                Clear();
                break;
            case 5:
                printf("\nTerima kasih telah menggunakan program antrian. Program selesai.\n");
                break;
            default:
                printf("[ERROR] Pilihan menu tidak valid! Silakan pilih antara 1 - 5.\n");
        }
    } while (pilihan != 5);

    return 0;
}
\end{lstlisting}

\vspace{0.4cm}
\noindent\textbf{2. Output Program:}

\begin{lstlisting}[style=cmdoutput]
========================================================
   PROGRAM IMPLEMENTASI QUEUE LINEAR ARRAY (SOAL 1)    
   Nama : Hellyos Ageng Haqiqie (NIM: 163251001)       
   Prodi: S1 Teknik Robotika dan Kecerdasan Buatan      
========================================================

========================================
           MENU PILIHAN UTAMA           
========================================
 1. Masukkan Data ke Antrian (Enqueue) 
 2. Hapus Data dari Antrian (Dequeue)  
 3. Tampilkan Data Antrian (Tampil)    
 4. Bersihkan Seluruh Antrian (Clear)  
 5. Keluar dari Program                
========================================
Pilihan Anda [1-5]: 1
Masukkan nilai bilangan bulat yang akan ditambahkan: 15
[SUKSES] Elemen 15 berhasil masuk ke antrian (posisi indeks: 0).

Pilihan Anda [1-5]: 1
Masukkan nilai bilangan bulat yang akan ditambahkan: 28
[SUKSES] Elemen 28 berhasil masuk ke antrian (posisi indeks: 1).

Pilihan Anda [1-5]: 1
Masukkan nilai bilangan bulat yang akan ditambahkan: 42
[SUKSES] Elemen 42 berhasil masuk ke antrian (posisi indeks: 2).

Pilihan Anda [1-5]: 1
Masukkan nilai bilangan bulat yang akan ditambahkan: 99
[SUKSES] Elemen 99 berhasil masuk ke antrian (posisi indeks: 3).

Pilihan Anda [1-5]: 1
Masukkan nilai bilangan bulat yang akan ditambahkan: 10
[SUKSES] Elemen 10 berhasil masuk ke antrian (posisi indeks: 4).

Pilihan Anda [1-5]: 3
=== STATUS ELEMEN ANTRIAN (LINEAR ARRAY) ===
Kapasitas Maksimum : 8 elemen
Jumlah Elemen Saat Ini: 5 elemen
Indeks Head (Depan)   : 0 (Nilai: 15)
Indeks Tail (Belakang): 4 (Nilai: 10)
Daftar Antrian: [FRONT] 15 -> 28 -> 42 -> 99 -> 10 [REAR]
Slot Memori Array:
+-------+-------+-------+-------+-------+-------+-------+-------+-------+
| Indeks|  [0]  |  [1]  |  [2]  |  [3]  |  [4]  |  [5]  |  [6]  |  [7]  |
+-------+-------+-------+-------+-------+-------+-------+-------+-------+
| Nilai |    15 |    28 |    42 |    99 |    10 |   -   |   -   |   -   |
+-------+-------+-------+-------+-------+-------+-------+-------+-------+

Pilihan Anda [1-5]: 2
[SUKSES] Elemen 15 berhasil keluar dari antrian.

Pilihan Anda [1-5]: 2
[SUKSES] Elemen 28 berhasil keluar dari antrian.

Pilihan Anda [1-5]: 3
=== STATUS ELEMEN ANTRIAN (LINEAR ARRAY) ===
Kapasitas Maksimum : 8 elemen
Jumlah Elemen Saat Ini: 3 elemen
Indeks Head (Depan)   : 0 (Nilai: 42)
Indeks Tail (Belakang): 2 (Nilai: 10)
Daftar Antrian: [FRONT] 42 -> 99 -> 10 [REAR]
Slot Memori Array:
+-------+-------+-------+-------+-------+-------+-------+-------+-------+
| Indeks|  [0]  |  [1]  |  [2]  |  [3]  |  [4]  |  [5]  |  [6]  |  [7]  |
+-------+-------+-------+-------+-------+-------+-------+-------+-------+
| Nilai |    42 |    99 |    10 |   -   |   -   |   -   |   -   |   -   |
+-------+-------+-------+-------+-------+-------+-------+-------+-------+

Pilihan Anda [1-5]: 4
[SUKSES] Seluruh isi antrian berhasil dibersihkan (Clear).

Pilihan Anda [1-5]: 3
[INFO] Antrian kosong! Tidak ada data untuk ditampilkan.

Pilihan Anda [1-5]: 5
Terima kasih telah menggunakan program antrian. Program selesai.
\end{lstlisting}

\newpage

% =============================================================================
% SOAL 2: OPERASI STATISTIK PADA ELEMEN QUEUE
% =============================================================================
\begin{center}
    {\Large\textbf{LATIHAN SOAL 2}}
\end{center}

\vspace{0.2cm}
\noindent\textbf{Soal:}
\begin{quote}
    \textit{"Carilah nilai total, rata-rata, terbesar dan terkecil dari elemen-elemen queue dengan menggunakan fungsi yang sudah dibuat pada soal 1!"}
\end{quote}

\vspace{0.3cm}
\noindent\textbf{1. Kode Program (\texttt{soal2\_queue\_statistik.c}):}

\begin{lstlisting}[style=ccode, language=C]
#include <stdio.h>
#include <stdlib.h>
#include <limits.h>

#define MAX 8

typedef struct {
    int data[MAX];
    int head;
    int tail;
} Queue;

Queue antrian;

void Create(void) {
    antrian.head = -1;
    antrian.tail = -1;
}

int IsEmpty(void) {
    return (antrian.tail == -1);
}

int IsFull(void) {
    return (antrian.tail >= MAX - 1);
}

void Enqueue(int data) {
    if (IsFull() == 1) {
        printf("[PERINGATAN] Antrian Penuh (Overflow)! Tidak dapat menambahkan data %d.\n", data);
        return;
    }
    if (IsEmpty() == 1) {
        antrian.head = 0;
        antrian.tail = 0;
        antrian.data[antrian.tail] = data;
    } else {
        antrian.tail++;
        antrian.data[antrian.tail] = data;
    }
    printf("[SUKSES] Elemen %d berhasil masuk ke antrian (indeks: %d).\n", data, antrian.tail);
}

int Dequeue(void) {
    if (IsEmpty() == 1) {
        printf("[PERINGATAN] Antrian Kosong (Underflow)!\n");
        return -1;
    }
    int elemenKeluar = antrian.data[antrian.head];
    for (int i = antrian.head; i < antrian.tail; i++) {
        antrian.data[i] = antrian.data[i + 1];
    }
    antrian.tail--;
    if (antrian.tail < 0) {
        antrian.head = -1;
        antrian.tail = -1;
    }
    printf("[SUKSES] Elemen %d berhasil dikeluarkan dari antrian.\n", elemenKeluar);
    return elemenKeluar;
}

void Clear(void) {
    antrian.head = -1;
    antrian.tail = -1;
    printf("[SUKSES] Seluruh isi antrian berhasil dibersihkan (Clear).\n");
}

void Tampil(void) {
    if (IsEmpty() == 1) {
        printf("[INFO] Antrian kosong!\n");
        return;
    }
    printf("\n=== STATUS ELEMEN ANTRIAN ===\n");
    printf("Kapasitas Array       : %d elemen\n", MAX);
    printf("Jumlah Elemen Aktif   : %d elemen\n", antrian.tail - antrian.head + 1);
    printf("Daftar Antrian (FIFO) : [FRONT] ");
    for (int i = antrian.head; i <= antrian.tail; i++) {
        printf("%d", antrian.data[i]);
        if (i < antrian.tail) printf(" -> ");
    }
    printf(" [REAR]\n");
}

/* FUNGSI STATISTIK (SOAL 2) */

/* Menghitung nilai total seluruh elemen di antrian */
long long HitungTotal(void) {
    if (IsEmpty() == 1) return 0;
    long long total = 0;
    for (int i = antrian.head; i <= antrian.tail; i++) {
        total += antrian.data[i];
    }
    return total;
}

/* Menghitung nilai rata-rata elemen di antrian */
double HitungRataRata(void) {
    if (IsEmpty() == 1) return 0.0;
    int jumlahElemen = antrian.tail - antrian.head + 1;
    return (double)HitungTotal() / (double)jumlahElemen;
}

/* Mencari nilai terbesar (maksimum) di antrian */
int CariTerbesar(void) {
    if (IsEmpty() == 1) return INT_MIN;
    int maxVal = antrian.data[antrian.head];
    for (int i = antrian.head + 1; i <= antrian.tail; i++) {
        if (antrian.data[i] > maxVal) {
            maxVal = antrian.data[i];
        }
    }
    return maxVal;
}

/* Mencari nilai terkecil (minimum) di antrian */
int CariTerkecil(void) {
    if (IsEmpty() == 1) return INT_MAX;
    int minVal = antrian.data[antrian.head];
    for (int i = antrian.head + 1; i <= antrian.tail; i++) {
        if (antrian.data[i] < minVal) {
            minVal = antrian.data[i];
        }
    }
    return minVal;
}

void TampilStatistik(void) {
    if (IsEmpty() == 1) {
        printf("\n[INFO] Antrian kosong! Tidak ada elemen untuk dianalisis.\n");
        return;
    }
    int jumlahElemen = antrian.tail - antrian.head + 1;
    long long total = HitungTotal();
    double rataRata = HitungRataRata();
    int terbesar = CariTerbesar();
    int terkecil = CariTerkecil();

    printf("\n=======================================================\n");
    printf("        HASIL ANALISIS STATISTIK ELEMEN ANTRIAN        \n");
    printf("=======================================================\n");
    printf("  Jumlah Elemen (Count)     : %d elemen\n", jumlahElemen);
    printf("  Nilai Total (Sum)         : %lld\n", total);
    printf("  Nilai Rata-rata (Mean)    : %.4f\n", rataRata);
    printf("  Nilai Terbesar (Maximum)  : %d\n", terbesar);
    printf("  Nilai Terkecil (Minimum)  : %d\n", terkecil);
    printf("=======================================================\n\n");
}

int main(void) {
    int pilihan, nilai;
    Create();

    printf("========================================================\n");
    printf("   PROGRAM QUEUE DENGAN OPERASI STATISTIK (SOAL 2)     \n");
    printf("   Nama : Hellyos Ageng Haqiqie (NIM: 163251001)       \n");
    printf("   Prodi: S1 Teknik Robotika dan Kecerdasan Buatan      \n");
    printf("========================================================\n");

    do {
        printf("\n========================================================\n");
        printf("                       MENU UTAMA                       \n");
        printf("========================================================\n");
        printf(" 1. Masukkan Data ke Antrian (Enqueue)                 \n");
        printf(" 2. Hapus Data dari Antrian (Dequeue)                  \n");
        printf(" 3. Tampilkan Seluruh Data Antrian (Tampil)            \n");
        printf(" 4. Hitung & Tampilkan Analisis Statistik (Total/Rata) \n");
        printf(" 5. Cari Nilai Terbesar Saja (Maksimum)                \n");
        printf(" 6. Cari Nilai Terkecil Saja (Minimum)                 \n");
        printf(" 7. Bersihkan Seluruh Antrian (Clear)                  \n");
        printf(" 8. Keluar dari Program                                \n");
        printf("========================================================\n");
        printf("Pilihan Anda [1-8]: ");
        if (scanf("%d", &pilihan) != 1) {
            printf("[ERROR] Input tidak valid!\n");
            while (getchar() != '\n');
            continue;
        }

        switch (pilihan) {
            case 1:
                printf("Masukkan bilangan bulat untuk ditambahkan: ");
                if (scanf("%d", &nilai) == 1) {
                    Enqueue(nilai);
                } else {
                    printf("[ERROR] Input tidak valid!\n");
                    while (getchar() != '\n');
                }
                break;
            case 2:
                Dequeue();
                break;
            case 3:
                Tampil();
                break;
            case 4:
                TampilStatistik();
                break;
            case 5:
                if (!IsEmpty()) {
                    printf("\n[HASIL] Nilai Terbesar (Maksimum) dalam antrian: %d\n", CariTerbesar());
                } else {
                    printf("\n[INFO] Antrian kosong!\n");
                }
                break;
            case 6:
                if (!IsEmpty()) {
                    printf("\n[HASIL] Nilai Terkecil (Minimum) dalam antrian: %d\n", CariTerkecil());
                } else {
                    printf("\n[INFO] Antrian kosong!\n");
                }
                break;
            case 7:
                Clear();
                break;
            case 8:
                printf("\nTerima kasih telah menggunakan program. Selesai.\n");
                break;
            default:
                printf("[ERROR] Pilihan menu tidak valid!\n");
        }
    } while (pilihan != 8);

    return 0;
}
\end{lstlisting}

\vspace{0.4cm}
\noindent\textbf{2. Output Program:}

\begin{lstlisting}[style=cmdoutput]
========================================================
   PROGRAM QUEUE DENGAN OPERASI STATISTIK (SOAL 2)     
   Nama : Hellyos Ageng Haqiqie (NIM: 163251001)       
   Prodi: S1 Teknik Robotika dan Kecerdasan Buatan      
========================================================

Pilihan Anda [1-8]: 1
Masukkan bilangan bulat untuk ditambahkan: 45
[SUKSES] Elemen 45 berhasil masuk ke antrian (indeks: 0).
Pilihan Anda [1-8]: 1
Masukkan bilangan bulat untuk ditambahkan: 12
[SUKSES] Elemen 12 berhasil masuk ke antrian (indeks: 1).
Pilihan Anda [1-8]: 1
Masukkan bilangan bulat untuk ditambahkan: 88
[SUKSES] Elemen 88 berhasil masuk ke antrian (indeks: 2).
Pilihan Anda [1-8]: 1
Masukkan bilangan bulat untuk ditambahkan: 34
[SUKSES] Elemen 34 berhasil masuk ke antrian (indeks: 3).
Pilihan Anda [1-8]: 1
Masukkan bilangan bulat untuk ditambahkan: 91
[SUKSES] Elemen 91 berhasil masuk ke antrian (indeks: 4).
Pilihan Anda [1-8]: 1
Masukkan bilangan bulat untuk ditambahkan: 23
[SUKSES] Elemen 23 berhasil masuk ke antrian (indeks: 5).
Pilihan Anda [1-8]: 1
Masukkan bilangan bulat untuk ditambahkan: 70
[SUKSES] Elemen 70 berhasil masuk ke antrian (indeks: 6).

Pilihan Anda [1-8]: 3
=== STATUS ELEMEN ANTRIAN ===
Kapasitas Array       : 8 elemen
Jumlah Elemen Aktif   : 7 elemen
Daftar Antrian (FIFO) : [FRONT] 45 -> 12 -> 88 -> 34 -> 91 -> 23 -> 70 [REAR]

Pilihan Anda [1-8]: 4
=======================================================
        HASIL ANALISIS STATISTIK ELEMEN ANTRIAN        
=======================================================
  Jumlah Elemen (Count)     : 7 elemen
  Nilai Total (Sum)         : 363
  Nilai Rata-rata (Mean)    : 51.8571
  Nilai Terbesar (Maximum)  : 91
  Nilai Terkecil (Minimum)  : 12
=======================================================

Pilihan Anda [1-8]: 5
[HASIL] Nilai Terbesar (Maksimum) dalam antrian: 91

Pilihan Anda [1-8]: 6
[HASIL] Nilai Terkecil (Minimum) dalam antrian: 12

Pilihan Anda [1-8]: 2
[SUKSES] Elemen 45 berhasil dikeluarkan dari antrian.

Pilihan Anda [1-8]: 3
=== STATUS ELEMEN ANTRIAN ===
Kapasitas Array       : 8 elemen
Jumlah Elemen Aktif   : 6 elemen
Daftar Antrian (FIFO) : [FRONT] 12 -> 88 -> 34 -> 91 -> 23 -> 70 [REAR]

Pilihan Anda [1-8]: 4
=======================================================
        HASIL ANALISIS STATISTIK ELEMEN ANTRIAN        
=======================================================
  Jumlah Elemen (Count)     : 6 elemen
  Nilai Total (Sum)         : 318
  Nilai Rata-rata (Mean)    : 53.0000
  Nilai Terbesar (Maximum)  : 91
  Nilai Terkecil (Minimum)  : 12
=======================================================

Pilihan Anda [1-8]: 8
Terima kasih telah menggunakan program. Selesai.
\end{lstlisting}

\newpage

% =============================================================================
% SOAL 3: CIRCULAR ARRAY DAN LINKED LIST QUEUE
% =============================================================================
\begin{center}
    {\Large\textbf{LATIHAN SOAL 3}}
\end{center}

\vspace{0.2cm}
\noindent\textbf{Soal:}
\begin{quote}
    \textit{"Buatlah program queue dengan menggunakan circular array dan linked list!"}
\end{quote}

\vspace{0.4cm}
\noindent{\large\textbf{3.1 Program Queue Menggunakan Circular Array}}

\vspace{0.2cm}
\noindent\textbf{1. Kode Program (\texttt{soal3\_circular\_queue.c}):}

\begin{lstlisting}[style=ccode, language=C]
#include <stdio.h>
#include <stdlib.h>

#define MAX 8

typedef struct {
    int data[MAX];
    int front; /* Menunjuk posisi sebelum elemen terdepan */
    int rear;  /* Menunjuk posisi elemen terakhir yang dimasukkan */
} CircularQueue;

CircularQueue q;

void Inisialisasi(void) {
    q.front = 0;
    q.rear = 0;
}

int IsEmpty(void) {
    return (q.front == q.rear);
}

int IsFull(void) {
    return ((q.rear + 1) % MAX == q.front);
}

int Noel(void) {
    return (q.rear - q.front + MAX) % MAX;
}

void Enqueue(int nilai) {
    if (IsFull()) {
        printf("[PERINGATAN] Circular Queue Penuh (OVERFLOW)! Kapasitas maksimum: %d elemen.\n", MAX - 1);
        printf("             Gagal menyisipkan nilai %d.\n", nilai);
        return;
    }
    q.rear = (q.rear + 1) % MAX;
    q.data[q.rear] = nilai;
    printf("[SUKSES] INSERT %d berhasil. Posisi Rear: [%d]. Total elemen = %d.\n", nilai, q.rear, Noel());
}

int Dequeue(void) {
    if (IsEmpty()) {
        printf("[PERINGATAN] Circular Queue Kosong (UNDERFLOW)!\n");
        return -1;
    }
    q.front = (q.front + 1) % MAX;
    int nilaiKeluar = q.data[q.front];
    printf("[SUKSES] REMOVE berhasil. Elemen %d dikeluarkan. Posisi Front: [%d]. Sisa = %d.\n",
           nilaiKeluar, q.front, Noel());
    return nilaiKeluar;
}

void Clear(void) {
    q.front = 0;
    q.rear = 0;
    printf("[SUKSES] Circular Queue berhasil dikosongkan (F = 0, R = 0).\n");
}

void Tampil(void) {
    if (IsEmpty()) {
        printf("[INFO] Circular Queue kosong (Front = %d, Rear = %d).\n", q.front, q.rear);
        return;
    }
    printf("\n=== STATUS CIRCULAR QUEUE ===\n");
    printf("Ukuran Array Fisik (MAX) : %d slot\n", MAX);
    printf("Kapasitas Maksimum Data  : %d elemen (MAX - 1)\n", MAX - 1);
    printf("Jumlah Elemen Aktif      : %d elemen\n", Noel());
    printf("Posisi Front             : %d\n", q.front);
    printf("Posisi Rear              : %d\n", q.rear);

    printf("Daftar Urutan Logis FIFO : [HEAD] ");
    int idx = (q.front + 1) % MAX;
    while (1) {
        printf("%d (slot [%d])", q.data[idx], idx);
        if (idx == q.rear) break;
        printf(" -> ");
        idx = (idx + 1) % MAX;
    }
    printf(" [TAIL]\n");

    printf("\nVisualisasi Memori Array Fisik:\n");
    printf("+-------");
    for (int i = 0; i < MAX; i++) printf("+-------");
    printf("+\n| Indeks");
    for (int i = 0; i < MAX; i++) printf("|  [%d]  ", i);
    printf("|\n+-------");
    for (int i = 0; i < MAX; i++) printf("+-------");
    printf("+\n| Nilai ");
    for (int i = 0; i < MAX; i++) {
        int isActive = 0;
        if (!IsEmpty()) {
            if (q.front < q.rear) {
                if (i > q.front && i <= q.rear) isActive = 1;
            } else {
                if (i > q.front || i <= q.rear) isActive = 1;
            }
        }
        if (isActive) printf("| %5d ", q.data[i]);
        else printf("|   -   ");
    }
    printf("|\n+-------");
    for (int i = 0; i < MAX; i++) printf("+-------");
    printf("+\n\n");
}

int main(void) {
    int pilihan, nilai;
    Inisialisasi();

    printf("========================================================\n");
    printf("   PROGRAM IMPLEMENTASI CIRCULAR QUEUE (SOAL 3A)       \n");
    printf("   Nama : Hellyos Ageng Haqiqie (NIM: 163251001)       \n");
    printf("   Prodi: S1 Teknik Robotika dan Kecerdasan Buatan      \n");
    printf("========================================================\n");

    do {
        printf("\n========================================================\n");
        printf("                 MENU PILIHAN CIRCULAR QUEUE            \n");
        printf("========================================================\n");
        printf(" 1. Sisipkan Elemen (Enqueue / INSERT)                 \n");
        printf(" 2. Ambil Elemen Terdepan (Dequeue / REMOVE)            \n");
        printf(" 3. Tampilkan Status & Visualisasi Array Melingkar     \n");
        printf(" 4. Cek Kondisi (IsEmpty / IsFull / Noel)              \n");
        printf(" 5. Bersihkan Antrian (Clear / RESET)                  \n");
        printf(" 6. Keluar                                             \n");
        printf("========================================================\n");
        printf("Pilihan Anda [1-6]: ");
        if (scanf("%d", &pilihan) != 1) {
            printf("[ERROR] Input tidak valid!\n");
            while (getchar() != '\n');
            continue;
        }

        switch (pilihan) {
            case 1:
                printf("Masukkan bilangan bulat untuk di-insert: ");
                if (scanf("%d", &nilai) == 1) Enqueue(nilai);
                else { while (getchar() != '\n'); }
                break;
            case 2:
                Dequeue();
                break;
            case 3:
                Tampil();
                break;
            case 4:
                printf("\n--- STATUS KONDISI ANTRIAN ---\n");
                printf("IsEmpty (Kosong) : %s\n", IsEmpty() ? "YA" : "TIDAK");
                printf("IsFull (Penuh)   : %s\n", IsFull() ? "YA" : "TIDAK");
                printf("Jumlah Elemen    : %d / %d\n", Noel(), MAX - 1);
                break;
            case 5:
                Clear();
                break;
            case 6:
                printf("\nProgram selesai.\n");
                break;
            default:
                printf("[ERROR] Pilihan menu tidak valid!\n");
        }
    } while (pilihan != 6);

    return 0;
}
\end{lstlisting}

\vspace{0.4cm}
\noindent\textbf{2. Output Program Circular Queue:}

\begin{lstlisting}[style=cmdoutput]
========================================================
   PROGRAM IMPLEMENTASI CIRCULAR QUEUE (SOAL 3A)       
   Nama : Hellyos Ageng Haqiqie (NIM: 163251001)       
   Prodi: S1 Teknik Robotika dan Kecerdasan Buatan      
========================================================

Pilihan Anda [1-6]: 1
Masukkan bilangan bulat untuk di-insert: 10
[SUKSES] INSERT 10 berhasil. Posisi Rear bergerak ke slot [1]. Total elemen = 1.
Pilihan Anda [1-6]: 1
Masukkan bilangan bulat untuk di-insert: 20
[SUKSES] INSERT 20 berhasil. Posisi Rear bergerak ke slot [2]. Total elemen = 2.
Pilihan Anda [1-6]: 1
Masukkan bilangan bulat untuk di-insert: 30
[SUKSES] INSERT 30 berhasil. Posisi Rear bergerak ke slot [3]. Total elemen = 3.
Pilihan Anda [1-6]: 1
Masukkan bilangan bulat untuk di-insert: 40
[SUKSES] INSERT 40 berhasil. Posisi Rear bergerak ke slot [4]. Total elemen = 4.
Pilihan Anda [1-6]: 1
Masukkan bilangan bulat untuk di-insert: 50
[SUKSES] INSERT 50 berhasil. Posisi Rear bergerak ke slot [5]. Total elemen = 5.
Pilihan Anda [1-6]: 1
Masukkan bilangan bulat untuk di-insert: 60
[SUKSES] INSERT 60 berhasil. Posisi Rear bergerak ke slot [6]. Total elemen = 6.
Pilihan Anda [1-6]: 1
Masukkan bilangan bulat untuk di-insert: 70
[SUKSES] INSERT 70 berhasil. Posisi Rear bergerak ke slot [7]. Total elemen = 7.

Pilihan Anda [1-6]: 3
=== STATUS CIRCULAR QUEUE ===
Ukuran Array Fisik (MAX) : 8 slot
Kapasitas Maksimum Data  : 7 elemen (MAX - 1)
Jumlah Elemen Aktif      : 7 elemen
Posisi Front             : 0
Posisi Rear              : 7
Daftar Urutan Logis FIFO : [HEAD] 10 (slot [1]) -> 20 (slot [2]) -> 30 (slot [3]) -> 40 (slot [4]) -> 50 (slot [5]) -> 60 (slot [6]) -> 70 (slot [7]) [TAIL]

Pilihan Anda [1-6]: 1
Masukkan bilangan bulat untuk di-insert: 999
[PERINGATAN] Circular Queue Penuh (OVERFLOW)! Kapasitas maksimum adalah 7 elemen.
             Gagal menyisipkan nilai 999.

Pilihan Anda [1-6]: 2
[SUKSES] REMOVE berhasil. Elemen 10 dikeluarkan. Posisi Front bergerak ke slot [1]. Sisa elemen = 6.
Pilihan Anda [1-6]: 2
[SUKSES] REMOVE berhasil. Elemen 20 dikeluarkan. Posisi Front bergerak ke slot [2]. Sisa elemen = 5.
Pilihan Anda [1-6]: 2
[SUKSES] REMOVE berhasil. Elemen 30 dikeluarkan. Posisi Front bergerak ke slot [3]. Sisa elemen = 4.

Pilihan Anda [1-6]: 1
Masukkan bilangan bulat untuk di-insert: 80
[SUKSES] INSERT 80 berhasil. Posisi Rear bergerak ke slot [0]. Total elemen = 5.

Pilihan Anda [1-6]: 1
Masukkan bilangan bulat untuk di-insert: 90
[SUKSES] INSERT 90 berhasil. Posisi Rear bergerak ke slot [1]. Total elemen = 6.

Pilihan Anda [1-6]: 3
=== STATUS CIRCULAR QUEUE ===
Ukuran Array Fisik (MAX) : 8 slot
Kapasitas Maksimum Data  : 7 elemen (MAX - 1)
Jumlah Elemen Aktif      : 6 elemen
Posisi Front             : 3
Posisi Rear              : 1
Daftar Urutan Logis FIFO : [HEAD] 40 (slot [4]) -> 50 (slot [5]) -> 60 (slot [6]) -> 70 (slot [7]) -> 80 (slot [0]) -> 90 (slot [1]) [TAIL]
\end{lstlisting}

\newpage

% -----------------------------------------------------------------------------
\noindent{\large\textbf{3.2 Program Queue Menggunakan Linked List}}

\vspace{0.2cm}
\noindent\textbf{1. Kode Program (\texttt{soal3\_linked\_list\_queue.c}):}

\begin{lstlisting}[style=ccode, language=C]
#include <stdio.h>
#include <stdlib.h>

typedef struct Node {
    int data;
    struct Node* next;
} Node;

typedef struct {
    Node* front;
    Node* rear;
    int count;
} LinkedQueue;

LinkedQueue antrian;

void InitQueue(LinkedQueue* q) {
    q->front = NULL;
    q->rear = NULL;
    q->count = 0;
}

int IsEmpty(LinkedQueue* q) {
    return (q->front == NULL);
}

void Enqueue(LinkedQueue* q, int data) {
    Node* newNode = (Node*)malloc(sizeof(Node));
    if (newNode == NULL) {
        printf("[ERROR] Alokasi memori gagal (Heap Overflow)!\n");
        return;
    }
    newNode->data = data;
    newNode->next = NULL;

    if (IsEmpty(q)) {
        q->front = newNode;
        q->rear = newNode;
    } else {
        q->rear->next = newNode;
        q->rear = newNode;
    }
    q->count++;
    printf("[SUKSES] Enqueue: Elemen %d ditambahkan pada alamat [%p]. Total: %d.\n",
           data, (void*)newNode, q->count);
}

int Dequeue(LinkedQueue* q) {
    if (IsEmpty(q)) {
        printf("[PERINGATAN] Antrian Kosong (Underflow)!\n");
        return -1;
    }
    Node* temp = q->front;
    int dataKeluar = temp->data;

    q->front = q->front->next;
    if (q->front == NULL) {
        q->rear = NULL;
    }
    free(temp);
    q->count--;
    printf("[SUKSES] Dequeue: Elemen %d berhasil dihapus. Sisa: %d.\n", dataKeluar, q->count);
    return dataKeluar;
}

int Peek(LinkedQueue* q) {
    if (IsEmpty(q)) {
        printf("[INFO] Antrian kosong!\n");
        return -1;
    }
    return q->front->data;
}

void Clear(LinkedQueue* q) {
    Node* curr = q->front;
    while (curr != NULL) {
        Node* nextNode = curr->next;
        free(curr);
        curr = nextNode;
    }
    q->front = NULL;
    q->rear = NULL;
    q->count = 0;
    printf("[SUKSES] Seluruh node memori berhasil didealokasikan (Clear).\n");
}

void Tampil(LinkedQueue* q) {
    if (IsEmpty(q)) {
        printf("[INFO] Antrian kosong! Tidak ada node.\n");
        return;
    }
    printf("\n=== STATUS LINKED LIST QUEUE ===\n");
    printf("Jumlah Elemen Aktif (Count) : %d node\n", q->count);
    printf("Alamat Node Front           : %p (Nilai: %d)\n", (void*)q->front, q->front->data);
    printf("Alamat Node Rear            : %p (Nilai: %d)\n", (void*)q->rear, q->rear->data);

    printf("\nAlur Rantai Antrian (FIFO):\n");
    printf("[FRONT] -> ");
    Node* curr = q->front;
    while (curr != NULL) {
        printf("[Data: %d | Next: %p]", curr->data, (void*)curr->next);
        if (curr->next != NULL) printf(" -> ");
        curr = curr->next;
    }
    printf(" -> NULL [REAR]\n\n");
}

int main(void) {
    int pilihan, nilai;
    InitQueue(&antrian);

    printf("========================================================\n");
    printf("   PROGRAM IMPLEMENTASI LINKED LIST QUEUE (SOAL 3B)    \n");
    printf("   Nama : Hellyos Ageng Haqiqie (NIM: 163251001)       \n");
    printf("   Prodi: S1 Teknik Robotika dan Kecerdasan Buatan      \n");
    printf("========================================================\n");

    do {
        printf("\n========================================================\n");
        printf("               MENU LINKED LIST QUEUE                   \n");
        printf("========================================================\n");
        printf(" 1. Tambah Elemen Baru (Enqueue)                       \n");
        printf(" 2. Ambil Elemen Terdepan (Dequeue)                    \n");
        printf(" 3. Lihat Elemen Terdepan (Peek)                       \n");
        printf(" 4. Tampilkan Seluruh Rantai Queue                     \n");
        printf(" 5. Kosongkan Seluruh Node (Clear)                     \n");
        printf(" 6. Keluar dari Program                                \n");
        printf("========================================================\n");
        printf("Pilihan Anda [1-6]: ");
        if (scanf("%d", &pilihan) != 1) {
            printf("[ERROR] Input tidak valid!\n");
            while (getchar() != '\n');
            continue;
        }

        switch (pilihan) {
            case 1:
                printf("Masukkan bilangan bulat untuk di-enqueue: ");
                if (scanf("%d", &nilai) == 1) Enqueue(&antrian, nilai);
                else { while (getchar() != '\n'); }
                break;
            case 2:
                Dequeue(&antrian);
                break;
            case 3:
                nilai = Peek(&antrian);
                if (nilai != -1) printf("\n[HASIL PEEK] Elemen terdepan: %d\n", nilai);
                break;
            case 4:
                Tampil(&antrian);
                break;
            case 5:
                Clear(&antrian);
                break;
            case 6:
                Clear(&antrian);
                printf("\nProgram selesai.\n");
                break;
            default:
                printf("[ERROR] Pilihan menu tidak valid!\n");
        }
    } while (pilihan != 6);

    return 0;
}
\end{lstlisting}

\vspace{0.4cm}
\noindent\textbf{2. Output Program Linked List Queue:}

\begin{lstlisting}[style=cmdoutput]
========================================================
   PROGRAM IMPLEMENTASI LINKED LIST QUEUE (SOAL 3B)    
   Nama : Hellyos Ageng Haqiqie (NIM: 163251001)       
   Prodi: S1 Teknik Robotika dan Kecerdasan Buatan      
========================================================

Pilihan Anda [1-6]: 1
Masukkan bilangan bulat untuk di-enqueue: 100
[SUKSES] Enqueue: Elemen 100 berhasil ditambahkan pada alamat memori [000002161A4F1470]. Total elemen: 1.
Pilihan Anda [1-6]: 1
Masukkan bilangan bulat untuk di-enqueue: 200
[SUKSES] Enqueue: Elemen 200 berhasil ditambahkan pada alamat memori [000002161A4F1490]. Total elemen: 2.
Pilihan Anda [1-6]: 1
Masukkan bilangan bulat untuk di-enqueue: 300
[SUKSES] Enqueue: Elemen 300 berhasil ditambahkan pada alamat memori [000002161A4F14B0]. Total elemen: 3.
Pilihan Anda [1-6]: 1
Masukkan bilangan bulat untuk di-enqueue: 400
[SUKSES] Enqueue: Elemen 400 berhasil ditambahkan pada alamat memori [000002161A4F14D0]. Total elemen: 4.

Pilihan Anda [1-6]: 4
=== STATUS LINKED LIST QUEUE ===
Jumlah Elemen Aktif (Count) : 4 node
Alamat Node Front           : 000002161A4F1470 (Nilai: 100)
Alamat Node Rear            : 000002161A4F14D0 (Nilai: 400)

Alur Rantai Antrian (FIFO):
[FRONT] -> [Data: 100 | Next: 000002161A4F1490] -> [Data: 200 | Next: 000002161A4F14B0] -> [Data: 300 | Next: 000002161A4F14D0] -> [Data: 400 | Next: 0000000000000000] -> NULL [REAR]

Pilihan Anda [1-6]: 3
[HASIL PEEK] Elemen terdepan adalah: 100

Pilihan Anda [1-6]: 2
[SUKSES] Dequeue: Elemen 100 berhasil dihapus dari depan antrian. Sisa elemen: 3.
Pilihan Anda [1-6]: 2
[SUKSES] Dequeue: Elemen 200 berhasil dihapus dari depan antrian. Sisa elemen: 2.

Pilihan Anda [1-6]: 1
Masukkan bilangan bulat untuk di-enqueue: 500
[SUKSES] Enqueue: Elemen 500 berhasil ditambahkan pada alamat memori [000002161A4F1470]. Total elemen: 3.
Pilihan Anda [1-6]: 1
Masukkan bilangan bulat untuk di-enqueue: 600
[SUKSES] Enqueue: Elemen 600 berhasil ditambahkan pada alamat memori [000002161A4F1490]. Total elemen: 4.

Pilihan Anda [1-6]: 4
=== STATUS LINKED LIST QUEUE ===
Jumlah Elemen Aktif (Count) : 4 node
Alamat Node Front           : 000002161A4F14B0 (Nilai: 300)
Alamat Node Rear            : 000002161A4F1490 (Nilai: 600)

Alur Rantai Antrian (FIFO):
[FRONT] -> [Data: 300 | Next: 000002161A4F14D0] -> [Data: 400 | Next: 000002161A4F1470] -> [Data: 500 | Next: 000002161A4F1490] -> [Data: 600 | Next: 0000000000000000] -> NULL [REAR]

Pilihan Anda [1-6]: 5
[SUKSES] Seluruh node memori berhasil didealokasikan (Queue Clear).

Pilihan Anda [1-6]: 6
Program selesai.
\end{lstlisting}

\end{document}
